% Part: lambda-calculus
% Chapter: representability
% Section: currying

\documentclass[../../../include/open-logic-section]{subfiles}

\begin{document}

\olfileid[pt]{lam}{rep}{cur} 
\olsection{Currying}

Uma abstração-$\lambd$ $\lambd[x][M]$ representa uma função de um
argumento, o que é bastante limitante quando queremos definir uma função que
aceita vários argumentos. Uma maneira de fazer isso seria estender o
cálculo-$\lambd$ para permitir a formação de pares, triplas etc.; nesse
caso, digamos, uma função ternária $\lambd[x][M]$ esperaria que seu argumento
fosse uma tripla. Contudo, é mais conveniente fazer isso por meio de
\emph{Currying}.

Consideremos um exemplo. Fingiremos por um momento que temos uma operação
$+$ no cálculo-$\lambd$. A função de adição é de $2$ argumentos, isto é, recebe dois
argumentos. Mas uma abstração-$\lambd$ só nos dá funções de um argumento: a
sintaxe não permite expressões como $\lambd[(x,y)][(x+y)]$. Podemos, porém,
considerar a função unária~$f_x(y)$ dada por $\lambd[y][(x+y)]$, que soma
$x$ a seu único argumento~$y$. Na verdade, essa não é uma única função, mas
uma família de funções diferentes, ``somar $x$'', uma para cada número~$x$.
Podemos agora definir outra função unária~$g$ como $\lambd[x][f_x]$.
Aplicada ao argumento $x$, $g(x)$ devolve a função~$f_x$---portanto, seus
valores são outras funções. Agora, se aplicarmos $g$ a $x$ e depois o
resultado a~$y$, obtemos: $(g(x))y = f_x(y) = x+y$. Desse modo, a função
unária~$g$ pode cumprir a mesma tarefa que a função binária de adição.
``Currying'' designa simplesmente esse artifício de transformar funções
binárias em funções unárias (cujos valores são funções unárias).

Eis um exemplo propriamente na sintaxe do cálculo-$\lambd$. Como
representamos a função $f(x,y) = x$? Se queremos definir uma função que
aceita dois argumentos e devolve o primeiro, podemos escrever
$\lambd[x][\lambd[y][x]]$, que literalmente é uma função que aceita um
argumento~$x$ e devolve a função~$\lambd[y][x]$. A função $\lambd[y][x]$
aceita outro argumento~$y$, mas o descarta e sempre devolve~$x$. Vejamos o
que acontece quando aplicamos $\lambd[x][\lambd[y][x]]$ a dois argumentos:
\begin{align*}
  (\lambd[x][\lambd[y][x]])MN 
  \bredone & (\lambd[y][M])N \\
  \bredone & M
\end{align*}

Em geral, para escrever uma função com parâmetros $x_1$, \dots,~$x_n$
definida por algum termo~$N$, podemos escrever
$\lambd[x_1][\lambd[x_2][\ldots\lambd[x_n][N]]]$. Se lhe aplicarmos $n$
argumentos, obtemos:
\begin{multline*}
  (\lambd[x_1][\lambd[x_2][\ldots\lambd[x_n][N]]]) M_1 \dots M_n \bredone\\
  \begin{aligned}
  \bredone {} & (\Subst{(\lambd[x_2][\ldots\lambd[x_n][N]])}{M_1}{x_1}) M_2
  \dots M_n\\
   \eqs {} & (\lambd[x_2][\ldots\lambd[x_n][\Subst{N}{M_1}{x_1}]]) M_2
            \dots M_n \\
  \vdots & \\
  \bredone {} &\Subst{\Subst{P}{M_1}{x_1}\ldots}{M_n}{x_n}
  \end{aligned}
\end{multline*}
A última linha significa literalmente tomar $M_i$ como substituto de $x_i$ no corpo
da definição da função, exatamente o que queremos ao aplicar vários
argumentos a uma função.
\end{document}

